% ==============================================================================
% Apêndice B: Planejamento do projeto
% ==============================================================================
\chapter{Planejamento do projeto}
\label{ap:planejamento}

\begin{instrucao}
\orientacao{Planejamento, viabilidade e riscos}
As competências do TCC incluem desenvolver um cronograma, rastrear o progresso
e gerenciar os riscos do projeto. Este apêndice mostra uma forma de registrar
esses elementos. Mantenha o cronograma realizado (e não apenas o planejado) e
comente no texto os desvios relevantes. Apêndices são material elaborado pelo
próprio autor e devem ser citados no texto pelo menos uma vez. Material de
outros autores vai em anexo.

Pela NBR 14724:2024, se um apêndice citar obras, as referências dessas obras
ficam no próprio apêndice, e não na lista geral de Referências.
\end{instrucao}

% ------------------------------------------------------------------------------
\section{Cronograma}

O \cref{qua:cronograma} apresenta o cronograma realizado, organizado nos seis
ciclos de desenvolvimento descritos na \cref{sec:introducao:metodo}. A
avaliação com usuários, prevista para o quinto ciclo, foi adiada para o sexto
porque a implantação em homologação levou mais tempo que o planejado.

\newcommand{\feito}{\cellcolor{gray!45}}
\begin{quadro}[htbp]
  \centering
  \caption{Cronograma realizado por ciclo de desenvolvimento}
  \label{qua:cronograma}
  \small
  \begin{tabular}{|l|c|c|c|c|c|c|}
    \hline
    \textbf{Atividade} & \textbf{C1} & \textbf{C2} & \textbf{C3} &
      \textbf{C4} & \textbf{C5} & \textbf{C6} \\
    \hline
    Levantamento de requisitos         & \feito & \feito &        &        &        &        \\
    \hline
    Protótipos e validação             &        & \feito &        &        &        &        \\
    \hline
    Projeto da arquitetura e dos dados &        & \feito & \feito &        &        &        \\
    \hline
    Implementação (RF01 a RF05)        &        &        & \feito & \feito &        &        \\
    \hline
    Implementação (RF06 a RF08)        &        &        &        & \feito & \feito &        \\
    \hline
    Testes automatizados               &        &        & \feito & \feito & \feito &        \\
    \hline
    Implantação em homologação         &        &        &        &        & \feito &        \\
    \hline
    Avaliação com usuários             &        &        &        &        &        & \feito \\
    \hline
    Escrita do texto                   & \feito &        &        & \feito & \feito & \feito \\
    \hline
  \end{tabular}
  \fonte{Elaborado pelo próprio autor. Cada ciclo teve duas semanas.}
\end{quadro}

% ------------------------------------------------------------------------------
\section{Análise de riscos}

O \cref{qua:riscos} apresenta os principais riscos identificados no início do
projeto, com a probabilidade, o impacto e as ações previstas. O risco R3 se
concretizou no quinto ciclo e a ação de contingência foi aplicada.

\begin{quadro}[htbp]
  \centering
  \caption{Riscos do projeto e ações de prevenção e contingência}
  \label{qua:riscos}
  \small
  \begin{tabularx}{\textwidth}{|c|>{\raggedright\arraybackslash}X|c|c|>{\raggedright\arraybackslash}X|}
    \hline
    \textbf{Código} & \textbf{Risco} & \textbf{Prob.} & \textbf{Impacto} &
      \textbf{Prevenção e contingência} \\
    \hline
    R1 & Pouca disponibilidade dos usuários para entrevistas & Média & Alto &
      Agendar as entrevistas no início do semestre. Usar questionário
      \emph{on-line} se não houver agenda. \\
    \hline
    R2 & Estimativas otimistas de esforço & Alta & Médio &
      Priorizar as histórias com os usuários. Adiar as de baixa prioridade. \\
    \hline
    R3 & Atraso na liberação do servidor de homologação & Média & Alto &
      Solicitar o servidor no segundo ciclo. Implantar em serviço gratuito de
      nuvem como alternativa temporária. \\
    \hline
    R4 & Falta de familiaridade com o arcabouço & Baixa & Médio &
      Fazer um protótipo técnico no primeiro ciclo. \\
    \hline
  \end{tabularx}
  \fonteautor
\end{quadro}
